\subsection{Tempered distributions}

DEFINITION 4. --- We call the space of tempered distributions on \(\mathbf{R}^{n}\) the dual space of \(\mathcal{S}(\mathbf{R})\) endowed with the topology of bounded convergence. It is denoted \(\mathcal{S}'(\mathbf{R}^{n})\) .

Since \(\mathcal{S}(\mathbf{R}^n)\) is bornological, the space \(\mathcal{S}'(\mathbf{R}^n)\) is complete (TVS, III, p. 24, cor. 1). Since \(\mathcal{S}(\mathbf{R}^n)\) is a Montel space, the same holds for \(\mathcal{S}'(\mathbf{R}^n)\) (EVT, IV, p. 19, prop. 9). The space \(\mathcal{S}'(\mathbf{R}^n)\) is therefore reflexive (EVT, IV, p. 16, th. 2).

Lemma 9. --- A linear map f from \(\mathcal{S}(\mathbf{R}^{n})\) into \(\mathbf{C}\) is a tempered distribution if and only if for every family \((\mathrm{M}_{\alpha,k})_{(\alpha,k)\in\mathbf{N}^{n}\times\mathbf{N}}\) in \(\mathbf{R}_{+}\) the linear form f is bounded on the set of functions \(\varphi\in\mathcal{S}(\mathbf{R}^{n})\) such that \(\widetilde{q}_{\alpha,k}(\varphi)\leqslant\mathrm{M}_{\alpha,k}\) for all \((\alpha,k)\in\mathbf{N}^{n}\times\mathbf{N}\) .

Indeed, the space \(\mathcal{S}(\mathbf{R}^{n})\) is bornological (EVT, III, p. 12, prop. 2) and every bounded subset of \(\mathcal{S}(\mathbf{R}^{n})\) is contained in one of the bounded sets described in the statement.

Lemma 10. --- Let F be a filter on \(\mathcal{S}'(\mathbf{R}^n)\) having a countable base or containing a simply bounded set. Then F converges to a tempered distribution if and only if \(\langle f, \varphi \rangle\) converges according to F for every Schwartz function \(\varphi\) .

In particular, a sequence \((f_{m})_{m\in\mathbf{N}}\) of tempered distributions converges to a tempered distribution f if and only if one has \(\langle f_{m},\varphi\rangle\to\langle f,\varphi\rangle\) for all \(\varphi\in\mathcal{S}(\mathbf{R}^{n})\) .

Since \(\mathcal{S}'(\mathbf{R}^n)\) is a Fréchet space, hence barrelled, (EVT, III, p. 25, cor. of prop. 2) the lemma follows from the Banach-Steinhaus theorem (EVT, III, p. 26, cor. 3 of th. 1).

The canonical injection \(j\) of \(\mathcal{D}(\mathbf{R}^{n})\) into \(\mathcal{S}(\mathbf{R}^{n})\) is continuous and its image is dense (Lemma 12 of IV, p. 212), hence the transpose of \(j\) , which is the restriction map of tempered distributions to the subspace \(\mathcal{D}(\mathbf{R}^{n})\) is an injective continuous linear map from \(\mathcal{S}'(\mathbf{R}^{n})\) into \(\mathcal{D}'(\mathbf{R}^{n})\) . We will identify \(\mathcal{S}'(\mathbf{R}^{n})\) with a subspace of \(\mathcal{D}'(\mathbf{R}^{n})\) via this map.

Let \(\alpha \in \mathbf{N}^n\) . The linear map \(\varphi \mapsto \partial^\alpha \varphi\) from \(\mathcal{S}(\mathbf{R}^n)\) into \(\mathcal{S}(\mathbf{R}^n)\) is continuous. Its transpose is therefore a continuous linear map from \(\mathcal{S}'(\mathbf{R}^n)\) into \(\mathcal{S}'(\mathbf{R}^n)\) (EVT, IV, p. 6, cor., \(b\) )). We denote \(\partial^\alpha\) the continuous linear map \((-1)^{|\alpha|}{}^{t}\partial^\alpha\) from \(\mathcal{S}'(\mathbf{R}^n)\) into \(\mathcal{S}'(\mathbf{R}^n)\) . This definition is compatible with definition 2 of IV, p. 208 for distributions.

Let \(h \in \mathcal{C}^{\infty}(\mathbf{R}^{n})\) be a function such that \(\partial^{\alpha}h\) is of polynomial growth for every \(\alpha \in \mathbf{N}^{n}\) . The transpose of the continuous linear map \(\varphi \mapsto h\varphi\) (prop. 14 of IV, p. 214) is a continuous linear map on \(\mathcal{S}'(\mathbf{R}^{n})\) , denoted \(f \mapsto hf\) .
% semantic-fragments: legacy-input-seam
\relax{}
